\documentclass[portrait,a1paper,fontscale=0.4]{baposter}

\usepackage{graphicx}
\usepackage{caption}

\graphicspath{{../../../logos},{figures/article2},{../../../articles/extended-abstract-article2/figures}}


\begin{document}
\begin{poster}%
	{%
		grid=false, columns=2, background=none, headershade=plain, textborder=none, headerborder=none, boxshade=none, eyecatcher=false}%
	{%
		\begin{minipage}{.15\textwidth}
			\center
			
\includegraphics[scale=0.3]{branding_eguk.gif.png}
			\vspace{3em}
			
\includegraphics[scale=0.3]{cgvc-logo.pdf}
		\end{minipage}
	}
	{%
		Model reevaluation based on graph transformation rules
	}
	{
		\vspace{1em}
		M.Gaide, D.Marcheix, A.Arnould, X.Skapin, H.Belhaouari \& S.Jean
	}
	% {}
	{
		\begin{minipage}{.2\textwidth}
			
\includegraphics[scale=0.35]{Logo_ISAE-ENSMA}
			\hfill
			\includegraphics[scale=0.11]{Logo-up}
			\vspace{1em}
			
\includegraphics[scale=0.07]{lias-logo}
			\hfill
			
\includegraphics[scale=0.12]{xlim-logo}
		\end{minipage}
	}

	\begin{posterbox}[name=abstract,column=0,row=0, span=2]{Abstract}
		In this paper, we extend the scope of naming problem studies to encompass
		rule-based graph transformation modelling systems. We propose a novel
		persistent naming method that capitalizes on the formalized operations of
		generalized maps and graph transformation rules. It enables a unique and
		homogeneous characterisation of entities across all dimensions.
	\end{posterbox}

	\begin{posterbox}[name=context, below=abstract,span=2]{Context}
		\begin{minipage}[t]{.49\linewidth}
			\center
			\includegraphics[scale=0.5]{persistent-naming-pb-1}
			\captionof{figure}{Initial Evaluation}
			\includegraphics[scale=0.5]{persistent-naming-pb-2}
			\captionof{figure}{Reevaluation}
		\end{minipage}
	\end{posterbox}

	\begin{posterbox}[name=evaluation, below=context, row=3, column=0]{Evaluation}
		\begin{center}
			\includegraphics[width=0.19\textwidth]{fil-rouge-short-create-face.pdf}
			% \captionof{figure}{1-create(pos)}
			\includegraphics[width=0.19\textwidth]{fil-rouge-short-extrude-face.pdf}
			% \captionof{figure}{2-extrude($PN_{1}$, vec)}
			\includegraphics[width=0.19\textwidth]{fil-rouge-short-triangulate-face.pdf}
			% \captionof{figure}{3-triangulation($PN_{2}$)}
			\includegraphics[width=0.19\textwidth]{fil-rouge-short-color-face.pdf}
			% \captionof{figure}{4-color($PN_{3}$)}

			\captionof{figure}{Evaluation. 1-create(pos); 2-extrude($PN_1$, vec);
				3-triangulate($PN_2$); 4-color($PN_3$)}

			\vspace{1em}

			\includegraphics[height=8.5cm]{pn7-short}

			\captionof{figure}{Evaluation DAG of the persistent name $PN_{3}$}
		\end{center}
	\end{posterbox}

	\begin{posterbox}[name=reevaluation, below=context, column=1]{Reevaluation}
		\begin{center}
			\includegraphics[width=0.19\textwidth]{fil-rouge-short-create-face.pdf}
			\includegraphics[width=0.19\textwidth]{fil-rouge-short-reeval-insert-vertex.pdf}
			\includegraphics[width=0.19\textwidth]{fil-rouge-short-reeval-extrude-face.pdf}
			\includegraphics[width=0.19\textwidth]{fil-rouge-short-reeval-triangulate-face.pdf}
			\includegraphics[width=0.19\textwidth]{fil-rouge-short-reeval-color-face.pdf}

			\captionof{figure}{Reevaluation. 1-create(pos); ADD1-insert(3);
				extrude($PN_1$, vec); triangulate($PN_2$); 4-color($PN_3$)}

			\vspace{1em}

			\includegraphics[height=8.5cm]{pn7-reval-short}
			\captionof{figure}{Reevaluation DAG of persistent name $PN_{3}$. }
		\end{center}
	\end{posterbox}

	\begin{posterbox}[name=conclusion,above=bottom]{Conclusion}
		\begin{itemize}
			\item Topology naming extension to rule based modelling systems
			\item Implementation of persistent naming scheme
			      \begin{itemize}
				      \item Unique identification of any topological cell
			      \end{itemize}
			\item Possibility to adapt modelling strategies
			      \begin{itemize}
				      \item Add topological parameters
				      \item Remove operations
			      \end{itemize}
		\end{itemize}
	\end{posterbox}

\end{poster}
\end{document}
